\section{Conclusion}

MA-ROPE addresses the challenge of preserving deadline-compliant task delivery in dense vehicular edge computing (VEC) networks by jointly optimizing RAT selection, partial computation offloading, shared RSU queue contention, and federated cooperative learning. 

By evaluating policies across multiple fleet densities, we demonstrate that single-RAT baseline policies suffer catastrophic latency collapse as shared MEC queues saturate. To quantify this, we introduced the Deadline Retention Ratio (DRR). MA-ROPE achieves a DRR near 1.0, preserving nearly all of its reference-density performance under dense contention, whereas baseline methods lose substantial efficacy. Independent ns-3 co-simulation confirms that these results are not artifacts of the analytical model.

While MA-ROPE increases the per-attempt link failure rate slightly by opportunistically using less reliable RATs, its fleet-wide failure rate is substantially lower due to fewer transmission attempts. The joint decision space adds no meaningful runtime penalty.

Several limitations remain for future work. The energy model and centralized MADDPG upper bounds require further validation and are left for future research. The N=8 ablation study in Section II serves as a historical reference, predating our final physics corrections. Future extensions should incorporate trace-driven mobility, non-IID task distributions, RSU handovers, and formal differential privacy guarantees beyond FedAvg.

Ultimately, heterogeneous RAT selection, partial offloading, and shared MEC contention must be optimized jointly to satisfy strict V2X deadlines. Density-aware scaling metrics like DRR provide a direct and essential measure of a policy's real-world viability.
